
\documentclass[preprint]{IEIE}
\usepackage{graphicx}
\usepackage{amsmath}
\usepackage{booktabs}
\usepackage{array}
\usepackage{url}

\title{우연가설을 압도하는 관내사전투표\\동일 득표쌍 반복과 원자료 감사 요구}
\author{김준호 \\ 소속 없음 \\ e-mail : junhokim0331@gmail.com}
\engtitle{Repeated Identical Vote Pairs Overwhelming the Chance Hypothesis\\and the Need for Source-Record Audit}
\engauthor{Junho Kim \\ No Affiliation}
\abstract{
본 논문은 2026년 제9회 전국동시지방선거 관내사전투표에서 관찰된 동일 득표쌍 반복이 과거 공식 개표자료의 경험적 기준선과 양립하는지 검토한다. 중앙선거관리위원회 선거통계시스템 공식 화면 HTML에서 12개 사건행과 6개 동일 득표쌍을 재현했고, 2014년부터 2025년까지 공개자료 81,701개 개표단위 행을 파싱해 과거 기준선을 구성했다. 2014·2018·2022년 시·도지사 선거 51개 선거구에서 한 선거구 안 동일 득표쌍의 최대 반복은 3쌍이었으며, 광주전남 5쌍 이상 반복의 포아송 근사 확률은 약 0.115\%로 추정된다. 과거 실제 득표쌍 풀의 비복원 재표본추출 200,000회에서도 5쌍 이상 반복은 0회였고, 95\% 상한은 약 0.0015\%였다. 또한 2020년과 2024년 국회의원선거에서는 분석 가능한 모든 지역구에서 더불어민주당 후보의 사전투표 양자득표율이 당일투표보다 높았다. 본 연구의 결론은 특정 부정행위의 법적 확정이 아니라, 공개자료만으로도 우연가설을 약화시키는 통계적 이상성이 존재하며 독립적 원자료 감사가 필요하다는 것이다.\\[2mm]\noindent 주요어: 관내사전투표, 동일 득표쌍, 선거무결성, 중복확률, 포아송 근사, 사전투표 편차, 원자료 감사
}

\begin{document}
\maketitle
\makeabstract

\section{문제 제기}

2026년 지방선거 이후 서로 다른 관내사전투표 개표단위에서 주요 후보 2명의 득표수가 동시에 동일하게 나타난 사례가 보도되었다. TV조선은 2026년 6월 8일 기사에서 인천과 전남 등 전국 12곳에서 동일 득표가 나타났고, 특히 광주전남에서는 10개 지역에서 민주당 민형배 후보와 국민의힘 이정현 후보가 같은 득표를 기록했다고 보도했다. 같은 기사에서 인천 송도1동과 송도2동은 민주당 박찬대 후보 3030표, 국민의힘 유정복 후보 1440표로 동일한 것으로 제시되었다.

본 연구의 관심은 정치적 판단이 아니라 통계적 판단이다. 질문은 다음과 같다.

\begin{enumerate}
\item 같은 선거, 같은 후보 조합, 같은 관내사전투표 유형에서 1·2위 후보 득표쌍이 5쌍 이상 반복되는 현상은 과거 자료에서 얼마나 자주 관찰되었는가?
\item 단순 중복확률 모형에서 이러한 반복은 어느 정도의 확률로 발생하는가?
\item 이 현상은 단순 우연으로 처리할 수 있는가, 아니면 원자료 검증을 요구하는 통계적 이상치인가?
\end{enumerate}

본 논문은 세 번째 질문에 대해 신중하게 답한다. 통계적 희박성은 부정행위의 충분조건이 아니다. 그러나 통계적 희박성이 충분히 강하면, 선거관리기관이 "우연"이라고만 설명하고 끝낼 수 없는 감사·검증 사유가 된다.

\subsection{연구가설}

본 연구의 1차 검정은 동일 득표쌍 반복에 관한 것이다.

\begin{itemize}
\item 귀무가설 \(H_0\): 2026년 광주전남 관내사전투표에서 관찰된 1·2위 후보 동일 득표쌍 5쌍은 과거 시·도지사 선거에서 관찰된 동일 득표쌍 발생 구조와 양립한다. 쉽게 말해, 과거 자료로 볼 때 충분히 나올 수 있는 사건이라는 가정이다.
\item 대립가설 \(H_1\): 2026년 광주전남 관내사전투표의 동일 득표쌍 5쌍은 과거 시·도지사 선거의 경험분포와 양립하기 어려운 이상치이다. 쉽게 말해, 과거 자료로는 설명하기 어려운 사건이라는 가정이다.
\end{itemize}

2차 검정은 사전투표와 선거일 당일투표의 득표율 차이에 관한 것이다.

\begin{itemize}
\item 귀무가설 \(H_0'\): 사전투표와 당일투표의 득표율 차이는 선거구별 유권자 선택과 일반적 투표행태 차이의 범위 안에서 설명된다.
\item 대립가설 \(H_1'\): 사전투표와 당일투표의 득표율 차이는 다수 선거구에서 같은 방향으로 반복되며, 추가 설명변수 또는 원자료 검증 없이는 충분히 설명되지 않는다.
\end{itemize}

본 논문은 \(H_1\)과 \(H_1'\)을 법적 결론이 아니라 감사 필요성을 판단하기 위한 통계적 명제로 다룬다.

\section{관측된 사건의 정의}

본 논문에서 분석하는 사건은 다음과 같이 정의한다.

\begin{itemize}
\item 선거: 2026년 제9회 전국동시지방선거
\item 투표유형: 관내사전투표
\item 비교 대상: 같은 광역단체장 선거 안의 읍면동 단위 개표 결과
\item 후보 기준: 해당 선거의 주요 1·2위 후보
\item 사건: 서로 다른 두 개표단위에서 1위 후보 득표수와 2위 후보 득표수가 동시에 동일함
\end{itemize}

공개 보도로 제기된 광주전남 사례는 다음 5쌍으로 정리된다. 본 연구는 이 값들을 중앙선거관리위원회 선거통계시스템의 개표단위별 개표결과 화면에서 다시 추출해 대조했다.

\begin{table*}[t]
\centering
\scriptsize
\resizebox{\linewidth}{!}{%
\begin{tabular}{p{0.190\linewidth}p{0.190\linewidth}p{0.190\linewidth}p{0.190\linewidth}p{0.190\linewidth}}
\toprule
쌍 & 개표단위 A & 개표단위 B & 민형배 & 이정현 \\
\midrule
1 & 광주 광산구 송정1동 & 고흥군 금산면 & 1401 & 120 \\
2 & 신안군 하의면 & 여수시 상일동 & 506 & 42 \\
3 & 화순군 이양면 & 강진군 병영면 & 444 & 46 \\
4 & 함평군 엄다면 & 장성군 북하면 & 606 & 57 \\
5 & 보성군 노동면 & 신안군 팔금면 & 356 & 42 \\
\bottomrule
\end{tabular}
}
\end{table*}

인천 송도 사례는 별도 선거의 1쌍이다.

\begin{table*}[t]
\centering
\scriptsize
\resizebox{\linewidth}{!}{%
\begin{tabular}{p{0.237\linewidth}p{0.237\linewidth}p{0.237\linewidth}p{0.237\linewidth}}
\toprule
개표단위 A & 개표단위 B & 박찬대 & 유정복 \\
\midrule
송도1동 & 송도2동 & 3030 & 1440 \\
\bottomrule
\end{tabular}
}
\end{table*}

따라서 본 논문은 두 수준을 분리한다.

\begin{itemize}
\item 광주전남 내부 사건: 같은 선거 안 5쌍
\item 전국 보도 사건: 서로 다른 선거를 포함한 6쌍
\end{itemize}

통계적으로 더 강한 분석 대상은 광주전남 내부 사건이다. 전국 6쌍은 여러 선거를 훑은 뒤 발견된 사례이므로, 별도 보정이 필요하다.

\section{자료와 재현 방법}

\subsection{과거 기준선 자료}

2014년부터 2025년까지 공개된 중앙선거관리위원회 및 공공데이터포털 선거 개표자료를 사용했다.

\begin{center}
\scriptsize
\begin{tabular}{>{\raggedright\arraybackslash}p{0.25\linewidth}>{\raggedright\arraybackslash}p{0.63\linewidth}}
\toprule
선거 & 분석 행 \\
\midrule
2014 지방 & 25,600 \\
2016 국회의원 & 3,661 \\
2017 대통령 & 3,684 \\
2018 지방 & 22,926 \\
2020 국회의원 & 3,681 \\
2022 대통령 & 3,703 \\
2022 지방 & 10,942 \\
2024 국회의원 & 3,754 \\
2025 대통령 & 3,750 \\
합계 & 81,701 \\
\bottomrule
\end{tabular}
\end{center}

분석에서 합계, 소계, 거소·선상, 국외·재외 투표는 제외했다. 관내사전투표와 관외사전투표를 분리했으며, 본 논문의 핵심 비교는 관내사전투표이다.

2026년 지방선거 공식 통합 XLSX 개표자료는 본 원고 작성 시점의 공공데이터포털 검색에서 기준선 자료와 같은 방식으로 확보되지 않았다. 그러나 중앙선거관리위원회 선거통계시스템의 "개표단위별 개표결과" 화면은 제9회 전국동시지방선거 결과를 제공하고 있었으며, 본 연구는 해당 화면의 POST 응답 HTML을 저장·파싱해 12개 사건값을 재확인했다. 이 절차는 \texttt{scripts/fetch\_nec\_2026\_duplicate\_cases.py}에 구현되어 있고, 산출물은 \texttt{outputs/nec\_2026\_reported\_duplicate\_cases.csv}와 \texttt{outputs/nec\_2026\_reported\_duplicate\_pairs.csv}이다. 따라서 2026년 사건값은 더 이상 보도 화면에만 의존하지 않는다. 다만 통합 XLSX 파일 또는 개표상황표 원본이 공개되면 같은 사건 정의로 추가 재검산해야 한다.

\subsection{추가 자료 포함·제외 기준}

표본 신뢰도를 높이기 위해 2014년 지방선거 개표자료를 추가로 편입했다. 2014년 지방선거는 전국 단위 사전투표가 적용된 초기 선거이며, 관내사전투표와 관외사전투표가 개표자료에 구분되어 있어 본 연구의 기준선에 포함할 수 있다. 2010년 지방선거와 2012년 대통령·국회의원선거는 공개자료가 존재하더라도 동일한 관내사전투표 제도 아래의 비교 대상이 아니므로 본 분석의 핵심 기준선에서는 제외했다.

이 선택은 연구결과에 유리한 방향의 표본 선택이 아니다. 실제로 2014년 지방선거를 포함하면 과거 시·도지사 동일 득표쌍 수가 증가하고, 2026년 사건의 추정 희박성은 완화된다. 그럼에도 2026년 광주전남 5쌍 반복은 확장 기준선의 과거 최대값을 초과한다.

\subsection{실제 1·2위 후보 기준}

초기 단순 분석에서는 파일에 기재된 앞 두 후보를 비교할 수 있다. 그러나 보도에서 말하는 "1·2위 후보"와 파일상 후보 기재 순서는 항상 같지 않다. 따라서 본 논문에서는 다음 절차를 사용했다.

\begin{enumerate}
\item 같은 선거, 같은 선거종류, 같은 선거구, 같은 투표유형으로 개표단위를 묶는다.
\item 해당 묶음 안에서 후보별 총득표를 계산한다.
\item 총득표 기준 실제 1위와 2위 후보를 선정한다.
\item 각 읍면동 관내사전투표 개표단위에서 이 두 후보의 득표쌍을 만든다.
\item 같은 득표쌍이 두 개 이상 개표단위에서 반복되는지 계산한다.
\end{enumerate}

이 기준은 보도된 사건과 가장 가깝다.

\subsection{주장과 검증자료의 대응}

본 연구의 각 주장은 다음 자료로 검증된다. 이 표를 두는 이유는 논증의 강도를 분리하기 위해서이다. 통계적으로 확인 가능한 명제와 행정·법적으로 추가 확인이 필요한 명제를 한 문장으로 섞으면, 반론자가 가장 약한 고리를 공격해 전체 결론을 무효화할 수 있다.

현재 패키지 안에서 재현되는 명제는 다음이다.

\begin{enumerate}
\item 2026년 12개 사건행은 중앙선거관리위원회 선거통계시스템 공식 화면값에서 재현된다.
\item 이 12개 사건행은 6개의 동일 득표쌍을 이룬다.
\item 과거 시·도지사 관내사전투표 51개 선거구에서 한 선거구 안 최대 반복은 3쌍이다.
\item 광주전남 5쌍 반복의 포아송 근사 꼬리확률은 \(N=393, K=100,944.8\) 기준 약 0.115\%이다.
\item 과거 시·도지사 실제 득표쌍 풀에서 393개 개표단위를 비복원 재표본추출한 200,000회 검정에서 5쌍 이상 반복은 한 번도 나타나지 않았다.
\item 2020·2024 총선에서 민주당 후보의 사전투표 양자득표율은 분석 가능한 모든 지역구에서 당일투표보다 높았다.
\end{enumerate}

반대로 아직 현재 자료만으로 확정되지 않는 명제는 "동일성이 자연적 개표 과정에서 발생했는지, 입력·표시 오류인지, 다른 원인인지"이다. 이 원인 명제는 개표상황표 원본, 1차 분류기 결과, 재확인표 배분 내역, 입력 로그가 있어야 판단할 수 있다. 각 명제와 재현 파일의 상세 대응은 제출 패키지의 \texttt{evidence\_matrix\_ko.md}에 따로 정리했다.

따라서 본 연구의 완결된 결론은 "원인 확정"이 아니라 "공식 화면값과 과거 기준선만으로도 원자료 감사가 필요한 통계적 이상치"이다. 이 결론은 현재 패키지 안의 자료로 재현되며, 원인 확정 명제는 아직 필요한 원자료가 공개되지 않았으므로 별도로 남겨 둔다.

\section{경험적 기준선}

\subsection{과거 시·도지사 선거}

2026년 광주전남 사례는 광역단체장 선거의 관내사전투표에서 관찰된 반복이므로, 가장 직접적인 비교 대상은 과거 시·도지사 선거이다.

2014년, 2018년, 2022년 지방선거의 시·도지사 관내사전투표 자료를 실제 1·2위 후보 기준으로 재분석한 결과는 다음과 같다.

\begin{center}
\scriptsize
\begin{tabular}{>{\raggedright\arraybackslash}p{0.25\linewidth}>{\raggedright\arraybackslash}p{0.63\linewidth}}
\toprule
항목 & 값 \\
\midrule
시·도지사 선거구 수 & 51 \\
총 비교쌍 수 & 1,514,172 \\
동일 득표쌍 수 & 15 \\
동일 득표쌍이 1개 이상 있는 선거구 & 9 \\
한 선거구 안 최대 동일 득표쌍 수 & 3 \\
한 선거구 안 5쌍 이상 & 0 \\
\bottomrule
\end{tabular}
\end{center}

관찰된 동일 득표쌍은 다음과 같다.

\begin{table*}[t]
\centering
\scriptsize
\resizebox{\linewidth}{!}{%
\begin{tabular}{p{0.158\linewidth}p{0.158\linewidth}p{0.158\linewidth}p{0.158\linewidth}p{0.158\linewidth}p{0.158\linewidth}}
\toprule
선거 & 선거구 & 1위 후보 & 2위 후보 & 동일 득표쌍 & 개표단위 \\
\midrule
2014 지방 & 강원도 & 최문순 & 최흥집 & 49, 80 & 두촌면, 김삿갓면 \\
2014 지방 & 강원도 & 최문순 & 최흥집 & 57, 95 & 부론면, 묵호동 \\
2014 지방 & 경상남도 & 홍준표 & 김경수 & 293, 83 & 부곡면, 웅양면 \\
2014 지방 & 경상북도 & 김관용 & 오중기 & 160, 23 & 초전면, 현동면 \\
2014 지방 & 경상북도 & 김관용 & 오중기 & 186, 28 & 지산동, 안덕면 \\
2014 지방 & 경상북도 & 김관용 & 오중기 & 342, 43 & 평화동, 봉양면 \\
2014 지방 & 전라남도 & 이낙연 & 이성수 & 124, 18 & 남면, 엄다면 \\
2014 지방 & 전라남도 & 이낙연 & 이성수 & 190, 33 & 북일면, 흑산면 \\
2014 지방 & 전라북도 & 송하진 & 박철곤 & 145, 41 & 웅포면, 보안면 \\
2014 지방 & 전라북도 & 송하진 & 박철곤 & 174, 27 & 이평면, 이백면 \\
2014 지방 & 전라북도 & 송하진 & 박철곤 & 174, 50 & 서수면, 중앙동 \\
2014 지방 & 충청남도 & 안희정 & 정진석 & 54, 61 & 시초면, 문산면 \\
2018 지방 & 대구광역시 & 권영진 & 임대윤 & 674, 559 & 관음동, 현풍면 \\
2022 지방 & 전라남도 & 김영록 & 이정현 & 377, 90 & 동복면, 장산면 \\
2022 지방 & 충청북도 & 김영환 & 노영민 & 319, 158 & 양산면, 내북면 \\
\bottomrule
\end{tabular}
}
\end{table*}

중요한 점은 과거 시·도지사 자료에서 동일 득표쌍이 아예 없지는 않았다는 것이다. 특히 2014년 지방선거를 추가하면 동일 득표쌍이 늘어난다. 그러나 한 선거구 안에서 5쌍이 반복된 사례는 관찰되지 않았다. 확장 기준선의 과거 최대치는 3쌍이다.

여기서 "과거에도 한 선거구 안 3쌍이 있었으니 5쌍도 자연스러운 것 아닌가"라는 반론이 가능하다. 그러나 이 반론은 관측값의 크기와 꼬리확률을 구분하지 못한다. 3쌍은 과거 51개 시·도지사 선거구 전체를 모았을 때 한 번 관찰된 최대치이고, 5쌍은 그 과거 최대치를 다시 두 단계 넘는 사건이다. 같은 기준선에서 \(N=393\), \(K=100,945\)를 적용하면 3쌍 이상은 약 4.23\%로 드물지만 관찰될 수 있는 범위에 있다. 반면 5쌍 이상은 약 0.115\%로, 약 37분의 1 수준까지 더 작아진다. 따라서 본 논문은 "3쌍도 이상하므로 5쌍도 이상하다"라고 주장하지 않는다. 오히려 반대로, 3쌍까지는 과거 자료 안에서 실제로 관찰된 상단으로 인정하고, 5쌍은 그 상단을 초과한 사건으로 분리해 다룬다.

\subsection{가능한 득표쌍 범위의 경험적 추정}

시·도지사 선거의 전체 비교쌍 수는 1,514,172개이고, 관찰된 동일 득표쌍은 15개이다. 단순화하여 가능한 1·2위 득표쌍의 실질적 범위를 \(K\)라고 두면, 기대되는 동일쌍 발생 수는 대략 다음과 같다.

\[
E(C) \approx \frac{N(N-1)}{2K}
\]

과거 시·도지사 자료에서 역산하면,

\[
\hat{K} \approx \frac{1,514,172}{15} \approx 100,945
\]

이 값은 엄밀한 모수라기보다 경험적 기준선이다. 실제 선거 데이터는 모든 지역이 같은 조건에서 독립적으로 나온 값이 아니며, 읍면동 규모와 지역 성향이 강하게 작용한다. 그럼에도 이 값은 "과거 시·도지사 선거에서 실제로 어느 정도 동일쌍이 발생했는가"를 반영한다.

\section{확률 계산}

\subsection{포아송 근사}

한 선거구 안에 \(N\)개의 관내사전 개표단위가 있고, 각 개표단위의 1·2위 득표쌍이 가능한 득표쌍 범위 \(K\) 안에서 나타난다고 하자. 서로 다른 두 개표단위가 같은 득표쌍을 갖는 횟수를 \(C\)라고 두면, 드문 중복 사건은 포아송 분포로 근사할 수 있다. 쉽게 말해, 많은 개표단위가 있을 때 같은 숫자쌍이 우연히 몇 번 겹칠 수 있는지를 계산하는 모형이다.

\[
C \sim \mathrm{Poisson}(\lambda)
\]

\[
\lambda = \frac{N(N-1)}{2K}
\]

관심 확률은 다음이다.

\[
P(C \geq 5)
\]

이는 광주전남 보도 사건, 즉 같은 선거 안 5쌍 이상의 동일 득표쌍이 나타날 확률에 해당한다.

\subsection{\(N=393\), \(K=100,945\) 기준}

여기서 \(N\)은 보도된 10개 사건행의 수가 아니라, 광주전남 시·도지사 관내사전투표에서 동일 득표쌍이 발생할 수 있었던 잠재 개표단위 수를 뜻한다. 즉 "이미 발견된 사례 수"가 아니라 "같은 선거구·같은 후보 조합·같은 관내사전투표 유형 안에서 서로 비교될 수 있었던 읍면동급 개표단위 수"이다.

본 연구는 이 \(N\)을 행정구역 근사가 아니라 중앙선거관리위원회 선거통계시스템 VCCP08 공식 HTML에서 직접 센 값으로 둔다. 광주광역시 5개 구와 전라남도 22개 시군의 시·도지사선거 개표단위별 개표결과를 모두 파싱한 뒤, \texttt{관내사전투표}이면서 \texttt{합계}가 아닌 행만 세었다. 그 결과 \texttt{outputs/nec\_2026\_gwangju\_jeonnam\_unit\_counts.csv}에서 광주 96개, 전남 297개, 합계 393개 관내사전 개표단위가 확인된다. 따라서 본문 계산의 기준값은 \(N=393\)이다.

\(K=100,945\)도 임의로 고른 값이 아니다. 앞 절의 과거 시·도지사 기준선에서 실제 1·2위 후보 득표쌍 비교쌍은 1,514,172개였고, 그 안에서 동일 득표쌍은 15개였다. 드문 충돌 모형에서는 기대 충돌 수가 대략 \(\binom{N}{2}/K\)이므로, 과거 자료의 실측 충돌률을 거꾸로 풀면

\[
\hat{K} \approx \frac{1,514,172}{15} \approx 100,945
\]

가 된다. 따라서 \(N=393\), \(K=100,945\)는 편의적으로 튀어나온 조합이 아니라, 선관위 공식 화면에서 센 광주전남 관내사전 개표단위 수와 과거 시·도지사 자료에서 역산한 경험적 \(K\)를 결합한 기준선이다.

\[
\lambda = \frac{393 \times 392}{2 \times 100,945} \approx 0.76307
\]

따라서

\[
P(C \geq 5) \approx 0.115\%
\]

\[
P(C \geq 6) \approx 0.0143\%
\]

직관적으로 바꾸면 다음과 같다.

\begin{center}
\scriptsize
\begin{tabular}{>{\raggedright\arraybackslash}p{0.293\linewidth}>{\raggedright\arraybackslash}p{0.293\linewidth}>{\raggedright\arraybackslash}p{0.293\linewidth}}
\toprule
사건 & 확률 & 역수 해석 \\
\midrule
5쌍 이상 & 약 0.115\% & 약 871번에 1번 \\
6쌍 이상 & 약 0.0143\% & 약 6,982번에 1번 \\
\bottomrule
\end{tabular}
\end{center}

이 계산에서 5쌍은 광주전남 내부 사건, 6쌍은 전국 보도 사건을 단순화한 참고값이다. 전국 6쌍은 여러 선거를 함께 훑은 뒤 발견된 값이므로 엄밀한 해석에는 별도 보정이 필요하다.

\subsection{정확 pair-collision 계산}

위 계산은 포아송 근사다. 따라서 "근사식이 희귀성을 과장한 것 아니냐"는 반론을 점검하기 위해, 같은 \(N=393\), 반올림 \(K=100,945\) 기준에서 정확 pair-collision 계산을 별도로 수행했다. 이 계산은 개표단위가 하나씩 배정될 때 아직 임계값에 도달하지 않은 득표쌍 점유 상태를 동적계획법으로 추적하고, 5쌍 이상에 도달하는 경로를 꼬리확률로 흡수한다.

\begin{center}
\scriptsize
\begin{tabular}{>{\raggedright\arraybackslash}p{0.293\linewidth}>{\raggedright\arraybackslash}p{0.293\linewidth}>{\raggedright\arraybackslash}p{0.293\linewidth}}
\toprule
계산 방식 & \(P(C \geq 5)\) & 해석 \\
\midrule
포아송 근사 & 0.0011484064 & 약 0.115\% \\
정확 pair-collision 계산 & 0.0012190884 & 약 0.122\% \\
\bottomrule
\end{tabular}
\end{center}

정확 계산값은 포아송 근사보다 약간 크다. 즉 본문의 포아송 근사는 5쌍 사건을 실제보다 더 희귀하게 보이도록 만든 계산이 아니다. 근사식을 exact 계산으로 바꾸어도 광주전남 5쌍 사건은 약 0.12\% 수준의 낮은 꼬리확률 사건으로 남는다.

\subsection{가능한 득표쌍 범위에 대한 민감도 분석}

가능한 득표쌍 범위 \(K\)를 어떻게 잡느냐에 따라 확률은 달라진다. \(N=393\)에서 \(K\)를 바꿔 계산하면 다음과 같다.

\begin{center}
\scriptsize
\begin{tabular}{>{\raggedright\arraybackslash}p{0.293\linewidth}>{\raggedright\arraybackslash}p{0.293\linewidth}>{\raggedright\arraybackslash}p{0.293\linewidth}}
\toprule
가능한 득표쌍 범위 \(K\) & \(P(C \geq 5)\) & \(P(C \geq 6)\) \\
\midrule
50,000 & 2.0550\% & 0.5056\% \\
100,000 & 0.1197\% & 0.0151\% \\
100,945 & 0.1148\% & 0.0143\% \\
200,000 & 0.0051\% & 0.00033\% \\
337,354 & 0.00043\% & 0.000016\% \\
500,000 & 0.000064\% & 0.0000016\% \\
\bottomrule
\end{tabular}
\end{center}

가장 보수적으로 \(K=50,000\)처럼 가능한 득표쌍 범위를 크게 줄여도 5쌍 이상은 약 2.05\%이다. 이는 불가능한 사건은 아니지만, 여전히 하위 2\% 안팎의 드문 사건이다. 2014년을 포함한 확장 시·도지사 자료에서 역산한 \(K=100,945\)를 적용하면 5쌍 이상은 약 0.115\%로 추정된다. 이는 2014년의 낮은 사전투표 규모까지 포함한 보수적 기준선에서도 감사 트리거로 보기 충분한 수준이다.

\subsection{개표단위 수 \(N\) 가정 변화 점검}

선관위 공식 HTML에서 확인한 기준값은 \(N=393\)이다. 다만 개표단위 수 가정이 결론을 얼마나 바꾸는지 보기 위해 \(K=100,945\)를 고정하고 \(N\)을 350에서 450까지 바꾼 민감도 분석도 추가한다.

\begin{table*}[t]
\centering
\scriptsize
\resizebox{\linewidth}{!}{%
\begin{tabular}{p{0.237\linewidth}p{0.237\linewidth}p{0.237\linewidth}p{0.237\linewidth}}
\toprule
관내사전 개표단위 수 \(N\) & \(\lambda\) & \(P(C \geq 5)\) & 역수 해석 \\
\midrule
350 & 0.6050 & 0.0410\% & 약 2,441번에 1번 \\
390 & 0.7515 & 0.1074\% & 약 931번에 1번 \\
393 & 0.7631 & 0.1148\% & 약 871번에 1번 \\
400 & 0.7905 & 0.1340\% & 약 746번에 1번 \\
430 & 0.9137 & 0.2500\% & 약 400번에 1번 \\
450 & 1.0008 & 0.3672\% & 약 272번에 1번 \\
\bottomrule
\end{tabular}
}
\end{table*}

이 표는 \(N\) 가정이 결과를 바꾸지만 결론의 방향을 뒤집지는 않음을 보여준다. \(N=450\)처럼 비교쌍 수를 크게 잡아도 5쌍 이상은 0.4\% 미만이다. 따라서 광주전남 5쌍 반복은 특정한 \(N=393\) 가정에만 의존하는 결론이 아니다.

\subsection{경험적 비복원 재표본추출 검정}

포아송 근사는 해석이 간단하지만, "가능한 득표쌍 범위 \(K\)"라는 모형 가정을 포함한다. 따라서 본 연구는 모형 의존성을 줄이기 위해 과거 시·도지사 관내사전투표의 실제 1·2위 득표쌍 풀을 직접 재표본추출했다.

절차는 다음과 같다.

\begin{enumerate}
\item 2014·2018·2022년 시·도지사 관내사전투표에서 실제 1·2위 후보 득표쌍 10,322개를 모은다.
\item 이 중 393개 개표단위를 비복원으로 무작위 추출한다.
\item 추출 표본 안에서 동일 득표쌍이 몇 개 그룹 반복되는지 센다.
\item 이 절차를 고정 난수시드로 200,000회 반복한다.
\end{enumerate}

비복원 추출을 사용한 이유는 같은 개표단위를 다시 뽑아 생기는 인위적 중복을 막기 위해서이다. 이 검정은 "과거 실제 시·도지사 관내사전투표 개표단위 중 393개를 뽑았을 때 5쌍 이상 반복이 얼마나 자주 나타나는가"를 직접 묻는다.

결과는 다음과 같다.

\begin{center}
\scriptsize
\begin{tabular}{>{\raggedright\arraybackslash}p{0.293\linewidth}>{\raggedright\arraybackslash}p{0.293\linewidth}>{\raggedright\arraybackslash}p{0.293\linewidth}}
\toprule
반복 기준 & 발생 횟수 / 200,000회 & 경험적 비율 \\
\midrule
3쌍 이상 & 77 & 0.0385\% \\
4쌍 이상 & 4 & 0.0020\% \\
5쌍 이상 & 0 & 관측해상도 기준 \(<0.0005\%\) \\
\bottomrule
\end{tabular}
\end{center}

5쌍 이상이 한 번도 관찰되지 않았다는 말은 실제 확률이 수학적으로 0이라는 뜻이 아니다. 이 재표본추출 설계에서 관찰 가능한 해상도는 \(1/200{,}000=0.0005\%\)이고, 보수적인 plus-one 추정은 \(1/200{,}001\), 즉 약 0.0005\%이다. 또한 "3의 법칙"을 적용하면 5쌍 이상 확률의 95\% 상한은 약 0.0015\%이다. 3의 법칙은 독립 시행에서 어떤 사건이 \(n\)번 중 한 번도 관찰되지 않았을 때, 그 사건의 95\% 상한을 대략 \(3/n\)으로 잡는 보수적 경험규칙이다. 이 결과는 포아송 근사보다 더 강하지만, 본 논문은 이를 단독 결론으로 사용하지 않는다. 핵심은 서로 다른 접근, 즉 과거 최대값 3쌍, 포아송 근사 0.115\%, 비복원 재표본추출 0/200,000회가 모두 광주전남 5쌍 반복을 흔한 사건으로 보지 않는다는 점이다.

\section{왜 "한 쌍"과 "다섯 쌍"은 다른가}

동일 득표쌍 논의에서 가장 흔한 오해는 "한 쌍도 나올 수 있으니 여러 쌍도 나올 수 있다"는 식의 선형 비교이다. 그러나 본 연구의 검정대상은 "동일 득표쌍이 하나라도 존재하는가"가 아니라 "같은 선거구, 같은 후보 조합, 같은 관내사전투표 유형 안에서 동일 득표쌍이 몇 번 반복되는가"이다. 전자는 넓은 공간에서 한 번 충돌하는 사건이고, 후자는 같은 공간 안에서 충돌이 여러 번 누적되는 꼬리 사건이다.

개표단위가 많으면 동일 득표쌍 한 쌍은 충분히 생길 수 있다. 서로 비교할 수 있는 개표단위 쌍의 수가 빠르게 증가하기 때문이다.

\[
\binom{N}{2} = \frac{N(N-1)}{2}
\]

예컨대 \(N=393\)이면 비교쌍은 77,028개이다. 이 정도 비교쌍에서는 어딘가 한 번 정도 같은 득표쌍이 나올 수 있다. 따라서 본 논문은 동일 득표쌍 1쌍의 존재 자체를 이상치로 보지 않는다. "한 쌍"은 감사의 출발점이 아니라, 대규모 비교에서 자연스럽게 생길 수 있는 배경사건이다.

공개 논의에서 자주 언급되는 2022년 대선 대구 서구 비산1동 제3·제4투표소 사례도 이 구분을 보여준다. 본 연구가 \texttt{data/pres2022.xlsx}를 검산한 결과 두 투표소는 실제로 이재명 131표, 윤석열 618표로 같았다. 그러나 이 사례는 대통령선거 선거일 투표구의 한 쌍 사례이고, 본 연구의 핵심 사건은 지방선거 시·도지사 관내사전투표 읍면동 개표단위에서 같은 후보 조합 5쌍이 반복된 사건이다. 따라서 비산1동 사례는 "한 쌍은 과거에도 있었다"는 반론을 정직하게 인정하게 해주지만, 광주전남 5쌍 검정에 합산하거나 그 결론을 무효화하는 근거는 아니다.

반면 "다섯 쌍"은 서로 다른 판단대상이다. 한 선거구 안 동일 득표쌍 수를 \(C\)라고 하면, 본 연구가 묻는 것은 \(C \ge 1\)이 아니라 \(C \ge 5\)이다. 포아송 근사에서 \(C\)의 기대값을 \(\lambda\)라고 할 때, \(m\)쌍 이상이 나올 확률은 다음의 꼬리확률이다.

\[
P(C \ge m) = 1 - e^{-\lambda}\sum_{c=0}^{m-1}\frac{\lambda^c}{c!}
\]

이 식에서 \(m\)이 1에서 3, 4, 5로 증가할 때 확률은 산술적으로 조금씩 줄어드는 것이 아니라 꼬리 방향으로 빠르게 작아진다. 본 연구의 확장 기준선 \(N=393\), \(K=100,945\), \(\lambda=0.76307\)에서는 다음과 같다.

\begin{table*}[t]
\centering
\scriptsize
\resizebox{\linewidth}{!}{%
\begin{tabular}{p{0.237\linewidth}p{0.237\linewidth}p{0.237\linewidth}p{0.237\linewidth}}
\toprule
기준 & 확률 & 직관적 빈도 & 해석 \\
\midrule
\(C \ge 1\) & 약 53.37\% & 약 1.9번에 1번 & 한 쌍은 충분히 가능 \\
\(C \ge 2\) & 약 17.88\% & 약 5.6번에 1번 & 드물지만 흔히 볼 수 있음 \\
\(C \ge 3\) & 약 4.23\% & 약 23.7번에 1번 & 과거자료의 관찰 상단 \\
\(C \ge 4\) & 약 0.773\% & 약 129번에 1번 & 강한 감사 신호 \\
\(C \ge 5\) & 약 0.115\% & 약 871번에 1번 & 과거 최대값 초과 사건 \\
\bottomrule
\end{tabular}
}
\end{table*}

따라서 3쌍과 5쌍의 차이는 "두 쌍 더 많다"가 아니다. 같은 기준선에서 3쌍 이상은 약 4.23\%이고 5쌍 이상은 약 0.115\%이므로, 5쌍 이상은 3쌍 이상보다 약 37배 더 드문 사건이다. 또한 2014년을 포함한 과거 시·도지사 51개 선거구의 경험분포에서 한 선거구 안 최대 반복은 3쌍이었고, 5쌍 이상은 관찰되지 않았다. 즉 3쌍은 "과거에 실제로 관찰된 자연적 상단"이고, 5쌍은 "그 상단을 초과한 새 사건"이다.

이 구분은 반론 대응에서도 중요하다. "과거에도 3쌍이 있었으니 5쌍도 괜찮다"는 주장은 기온으로 치면 "과거 최고 38도가 있었으니 40도도 같은 수준"이라고 말하는 것과 비슷하다. 38도와 40도는 숫자로는 2도 차이지만, 최고기록을 넘는 순간 통계적 해석은 달라진다. 마찬가지로 동일 득표쌍 반복에서도 3쌍은 기준선을 설정하는 관측값이고, 5쌍은 그 기준선을 검정하는 관측값이다.

또 하나의 차이는 사후탐색 문제의 크기이다. 전국 모든 선거와 모든 후보를 훑은 뒤 "어딘가 한 쌍"을 찾았다면 이는 발견기회가 매우 큰 사건이므로 강한 증거가 되기 어렵다. 그러나 광주전남 5쌍은 같은 후보 조합, 같은 관내사전투표 유형, 같은 시·도지사 선거구 안에 집중된 반복으로 정의된다. 분석범위를 이렇게 좁히면 단순히 "많이 뒤져서 하나 걸린 사건"이라는 반론이 약해진다. 본 논문의 핵심 주장이 전국 12곳 전체가 아니라 광주전남 내부 5쌍에 놓이는 이유가 여기에 있다.

결론적으로, 본 연구는 "한 쌍이 나왔으니 이상하다"고 말하지 않는다. 한 쌍은 가능하다. 3쌍도 과거자료에서 관찰된 드문 상단으로 인정한다. 문제는 5쌍이 같은 구조 안에 집중되면서 과거 최대 관측값과 꼬리확률 기준을 동시에 넘었다는 점이다. 이것이 "한 쌍"과 "다섯 쌍"을 같은 사건으로 취급할 수 없는 통계적 이유이다.

\section{인천 송도 사례의 특이성}

인천 송도1동과 송도2동 사례는 별도 선거이므로 광주전남 5쌍 사건과 직접 합쳐서는 안 된다. 그러나 검증 관점에서는 중요한 단서를 제공한다.

언론 보도에 따르면 송도1동과 송도2동의 최종 관내사전투표 결과는 박찬대 3030표, 유정복 1440표로 같았다. 그러나 선거인 수, 제3후보 득표수, 무효표, 기권표는 달랐다. 또한 인천시선관위가 공개한 설명에 따르면 1차 투표지분류기 결과는 서로 같지 않았고, 재확인 대상 투표지를 수개표·합산한 뒤 최종적으로 두 후보의 득표수가 같아졌다.

보도된 1차 결과는 다음과 같다.

\begin{table*}[t]
\centering
\scriptsize
\resizebox{\linewidth}{!}{%
\begin{tabular}{p{0.237\linewidth}p{0.237\linewidth}p{0.237\linewidth}p{0.237\linewidth}}
\toprule
개표단위 & 박찬대 1차 & 유정복 1차 & 재확인 대상 \\
\midrule
송도1동 & 3016 & 1427 & 44 \\
송도2동 & 3011 & 1429 & 53 \\
\bottomrule
\end{tabular}
}
\end{table*}

최종 결과는 두 곳 모두 다음과 같다.

\begin{center}
\scriptsize
\begin{tabular}{>{\raggedright\arraybackslash}p{0.293\linewidth}>{\raggedright\arraybackslash}p{0.293\linewidth}>{\raggedright\arraybackslash}p{0.293\linewidth}}
\toprule
개표단위 & 박찬대 최종 & 유정복 최종 \\
\midrule
송도1동 & 3030 & 1440 \\
송도2동 & 3030 & 1440 \\
\bottomrule
\end{tabular}
\end{center}

송도 사례의 통계적 특이성은 두 방식으로 볼 수 있다. 첫째, 연수구 안에서 "어딘가 한 쌍"의 동일 득표쌍이 생길 확률이다. 선관위 공식 HTML에서 연수구 시·도지사선거 관내사전 개표단위는 15개로 파싱된다. 같은 \(K=100,944.8\) 기준선을 적용하면 비교쌍 수는 \(\binom{15}{2}=105\)이고,

\[
\lambda_{\text{연수}} = \frac{105}{100,944.8} \approx 0.001040
\]

이다. 따라서 연수구 내부에서 적어도 한 쌍의 동일 득표쌍이 나올 확률은

\[
P(C \ge 1) \approx 1-e^{-0.001040} \approx 0.104\%
\]

로, 약 962번에 1번 수준이다.

둘째, 이미 특정된 두 개표단위인 송도1동과 송도2동이 같은 1·2위 득표쌍을 갖는 조건부 확률이다. 같은 \(K\) 기준선에서 한 쌍의 두 단위가 서로 같은 득표쌍을 가질 확률은

\[
\frac{1}{K} \approx \frac{1}{100,944.8} \approx 0.000991\%
\]

로, 약 100,945번에 1번 수준이다. 이 값은 "송도1동과 송도2동"이라는 두 단위가 사전에 관심 대상이었다고 볼 때의 계산이다. 반대로 전국 모든 지역을 훑은 뒤 발견한 사건으로 보면 사후탐색 보정이 필요하므로 이 숫자를 그대로 전국 단정에 쓰면 안 된다.

따라서 송도 사례는 광주전남 5쌍 검정과 같은 층위의 주검정은 아니지만, 보조적 감사 신호로는 강하다. 특히 두 동은 선거인 수, 투표수, 무효표, 기권표가 서로 다른데도 1·2위 후보 최종 득표쌍만 같았다. 이 때문에 송도 사례의 핵심은 "최종 숫자가 우연히 같을 수 있는가"를 넘어서, 서로 다른 1차 결과와 재확인 대상 투표지가 어떤 후보별 배분을 거쳐 동일한 최종쌍을 만들었는지에 있다.

이 설명은 두 가지 방향으로 해석될 수 있다. 선관위 관점에서는 "1차 결과와 재확인 대상 수가 다르므로 단순 복사나 전산 오류가 아니다"라는 해명이 된다. 반대로 검증 관점에서는 "최종 동일성이 재확인 단계에서 만들어졌으므로 재확인표 배분 내역을 공개해야 한다"는 요구가 된다.

즉 송도 사례의 핵심 검증자료는 최종 득표표가 아니라 재확인 대상 투표지의 후보별 배분 내역이다.

\section{부정행위의 증거와 통계적 이상치의 구분}

본 연구의 결과는 다음 명제를 지지한다.

\begin{quote}
광주전남 5쌍 동일 득표 사건은 우연가설 아래에서 매우 낮은 확률의 사건이다.
\end{quote}

그러나 다음 명제까지 자동으로 지지하지는 않는다.

\begin{quote}
광주전남 5쌍 동일 득표 사건은 곧바로 부정선거의 증거이다.
\end{quote}

따라서 본 논문에서 "증거"라는 표현은 두 층위로 나누어 사용해야 한다. 첫째, 동일 득표쌍 반복은 우연가설을 약화시키는 통계적 증거이다. 둘째, 특정인이 특정 방식으로 개표를 조작했다는 직접 증거는 아니다. 3쌍 반복과 5쌍 반복은 모두 첫 번째 의미에서는 검토 대상이지만, 강도는 다르다. 3쌍은 과거 자료 안에서 관찰된 드문 상단이고, 5쌍은 그 상단을 넘은 사건이다. 따라서 3쌍과 5쌍을 모두 법적 단정으로 부르면 과장이고, 둘을 모두 "아무 의미 없는 우연"이라고 부르면 통계적 정보를 버리는 것이다.

통계적 이상치는 원인에 대한 질문을 만든다. 가능한 원인은 여러 가지다.

\begin{enumerate}
\item 매우 낮은 확률의 우연
\item 개표상황표 작성 또는 전산 입력 오류
\item 공개 시스템 표시 오류
\item 재확인표 처리 과정의 체계적 쏠림
\item 실제 부정행위
\end{enumerate}

현재 공개 보도만으로는 1번부터 5번 중 하나를 확정할 수 없다. 그러나 1번, 즉 순수 우연만으로 설명하기에는 확률이 너무 낮다. 따라서 원자료 공개와 감사가 필요하다.

\section{검증에 필요한 원자료}

통계적 의혹을 법적·행정적 증거로 전환하려면 다음 자료가 필요하다.

\begin{enumerate}
\item 12곳 전체의 개표상황표 원본
\item 각 개표단위의 1차 투표지분류기 후보별 결과
\item 재확인 대상 투표지 수
\item 재확인 대상 투표지의 후보별 최종 배분 내역
\item 무효표, 기권표, 제3후보 득표수
\item 개표부서, 분류기, 심사·집계부, 입력자, 시간대 로그
\item 선관위 전산 공개값과 개표상황표 원본의 대조 결과
\item 정당·후보자 추천 참관인의 이의제기 여부
\end{enumerate}

특히 중요한 것은 1차 결과와 최종 결과의 차이다. 만약 광주전남 5쌍에서도 송도처럼 1차 결과는 서로 달랐지만 재확인 또는 최종 합산 과정에서 동일 득표쌍이 만들어졌다면, 이는 단순 우연보다 더 강한 감사 사유가 된다.

\section{한계}

본 논문은 다음 한계를 가진다.

첫째, 2026년 지방선거의 공공데이터포털 공식 통합 XLSX 개표파일은 아직 직접 파싱하지 못했다. 그러나 2026년 사례값은 중앙선거관리위원회 선거통계시스템의 개표단위별 개표결과 HTML을 직접 저장·파싱해 대조했다. 따라서 현재 한계는 "언론 보도값만 사용했다"가 아니라 "공식 화면값은 확인했지만 통합 XLSX 파일과 개표상황표 원본까지는 확보하지 못했다"는 것이다. 최종 행정·법적 판단에는 중앙선거관리위원회 공식 통합 파일, 개표상황표 원본, 1차 분류기 결과와의 대조가 추가로 필요하다.

다만 이 한계는 현 단계 분석을 폐기하는 사유라기보다 결론의 범위를 정하는 사유이다. 공식 화면값 기준으로 12개 사건행과 6개 동일쌍은 재현되며, 통합 XLSX 파일이 공개되면 본 논문의 사건 정의와 재현 스크립트를 그대로 적용해 동일 득표쌍 반복 수를 다시 확인할 수 있다.

둘째, 포아송 근사는 단순화된 모형이다. 실제 읍면동 득표쌍은 모든 지역이 같은 조건에서 따로따로 만들어진 값이 아니며, 지역 성향, 투표자 수, 후보 수, 도시·농촌 차이, 사전투표율의 영향을 받는다.

셋째, 가능한 득표쌍 범위 \(K\)는 직접 관측되는 자연상수가 아니다. 본 논문은 과거 시·도지사 자료에서 경험적으로 역산한 값을 사용했다. 보수적 민감도 분석을 함께 제시했지만, 최종 연구에서는 과거 자료를 반복 재표본추출하는 방법 또는 지역별 규모 차이를 반영하는 계층모형으로 보완해야 한다.

넷째, "전국 12곳"은 여러 선거를 훑은 뒤 발견된 사례이므로 나중에 찾아낸 사례라는 문제가 있다. 본 논문에서 통계적으로 가장 의미 있게 다룬 것은 광주전남 한 선거구 내부의 5쌍 반복이다.

\section{사전투표-당일투표 득표율 차이 검정}

동일 득표쌍 반복과 별도로, 본 연구는 사전투표와 선거일 당일투표의 득표율 차이를 보조 검정으로 사용한다. 이 검정의 목적은 사전투표가 당일투표와 같은 방향·같은 규모의 정치적 선택을 반영했는지를 확인하는 것이다. 어떤 선거구에서 특정 후보의 사전투표 득표율을 \(p_E\), 당일투표 득표율을 \(p_D\), 사전투표 수를 \(n_E\), 당일투표 수를 \(n_D\)라고 하자. 두 투표군이 같은 기저 지지율 \(p\)에서 나온 무작위 표본이라는 귀무가설을 세우면, 표준적인 두 비율 차이 z통계량은 다음과 같다.

\[
z = \frac{p_E - p_D}{\sqrt{\hat{p}(1-\hat{p})(1/n_E + 1/n_D)}}
\]

여기서 \(\hat{p}\)는 두 표본을 합친 후보 득표율이다. 표본 수가 수천 또는 수만 단위라면 분모는 매우 작아진다. 따라서 사전투표와 당일투표 득표율 차이가 몇 퍼센트포인트만 나도 z값은 매우 커질 수 있다. 이 때문에 "큰 표본에서 큰 득표율 차이"는 단순 표본오차로 설명하기 어렵다는 주장이 나온다.

그러나 이 검정에는 중요한 전제가 있다. 사전투표자와 당일투표자가 같은 유권자 집단에서 무작위로 갈라져야 한다. 실제 선거에서는 이 전제가 그대로 성립하지 않을 수 있다. 사전투표자는 연령, 직업, 지역 이동성, 정당 선호, 캠페인 동원 전략에 따라 스스로 사전투표를 선택할 수 있다. 따라서 사전투표와 당일투표의 득표율 차이가 크다는 사실만으로는 곧바로 부정행위를 추론할 수 없다.

그럼에도 이 검정이 완전히 무의미한 것은 아니다. 다음 조건이 함께 관찰되면 의혹의 강도는 높아진다.

\begin{enumerate}
\item 사전투표와 당일투표의 차이가 여러 선거구에서 같은 정당 방향으로 반복된다.
\item 차이가 과거 같은 지역의 사전-당일 차이보다 갑자기 커진다.
\item 연령, 성별, 지역, 도시·농촌, 사전투표율 등 설명 가능한 요인을 반영해도 남는 차이가 크다.
\item 관내사전투표와 관외사전투표의 패턴이 서로 다르게 나타난다.
\item 동일 득표쌍 반복, 투표용지 부족, 개표상황표 불일치 같은 별도 이상치와 같은 지역·같은 선거에서 겹친다.
\end{enumerate}

따라서 본 논문은 사전투표-당일투표 득표율 차이를 보조 검정으로 위치시킨다. 동일 득표쌍 반복은 정수 득표쌍의 중복 문제이고, 사전-당일 득표율 차이는 두 표본 비율 차이 문제이다. 두 현상은 통계적 구조가 다르다. 그러나 둘이 같은 선거에서 동시에 관찰되면, 단일한 우연 설명은 더 약해진다.

\subsection{공식 총선자료 검산}

본 연구는 2016·2020·2024 국회의원선거 공식 개표자료를 직접 파싱했다. 분석 대상은 더불어민주당 후보와 보수정당 후보가 모두 존재하는 지역구이다. 보수정당은 2016년 새누리당, 2020년 미래통합당, 2024년 국민의힘으로 정의했다.

각 지역구에서 다음 값을 계산했다.

\[
\Delta = p_E - p_D
\]

여기서 \(p_E\)는 사전투표에서 민주당 후보가 민주당+보수정당 양자구도 표 중 얻은 비율이고, \(p_D\)는 선거일 당일투표의 같은 비율이다. 사전투표는 관내사전투표와 관외사전투표를 합산했고, 당일투표는 선거일 투표구 투표만 합산했다. 거소·선상, 국외·재외, 합계·소계는 제외했다.

직접 검산 결과는 다음과 같다.

\begin{table*}[t]
\centering
\scriptsize
\resizebox{\linewidth}{!}{%
\begin{tabular}{p{0.106\linewidth}p{0.106\linewidth}p{0.106\linewidth}p{0.106\linewidth}p{0.106\linewidth}p{0.106\linewidth}p{0.106\linewidth}p{0.106\linewidth}p{0.106\linewidth}}
\toprule
선거 & 지역구 수 & 평균 사전-당일 차이 & 중앙값 & 민주당 사전 우위 & 민주당 사전 열위 & |z|>5 & |z|>10 & 최대 |z| \\
\midrule
2016 총선 & 229 & +3.90\%p & +3.43\%p & 211 & 18 & 158 & 88 & 31.94 \\
2020 총선 & 236 & +10.88\%p & +11.21\%p & 236 & 0 & 236 & 234 & 55.77 \\
2024 총선 & 245 & +10.35\%p & +10.80\%p & 245 & 0 & 245 & 245 & 60.92 \\
\bottomrule
\end{tabular}
}
\end{table*}

이 결과는 두 가지를 동시에 보여준다.

첫째, "2016년에는 대수의 법칙이 완전히 성립했고 2020년부터만 깨졌다"는 식의 단순화는 공식자료와 맞지 않는다. 2016년에도 민주당 후보의 사전투표 양자득표율은 당일투표보다 평균 3.90\%p 높았고, 229개 지역구 중 211개에서 민주당 사전 우위가 나타났다.

둘째, 2020년과 2024년의 패턴은 2016년보다 훨씬 강하다. 2020년에는 분석 가능한 236개 지역구 전부에서 민주당 사전 우위가 나타났고, 2024년에도 245개 지역구 전부에서 민주당 사전 우위가 나타났다. 평균 차이도 약 10\textasciitilde{}11\%p로 확대되었다.

따라서 직접 검산으로 방어 가능한 명제는 다음이다.

\begin{quote}
2016년에도 민주당 사전투표 우위는 존재했으나, 2020년과 2024년에는 그 방향성이 전 지역구로 확대되고 크기도 약 3배 커졌다.
\end{quote}

이 명제는 통계적으로 강하다. 다만 이 역시 곧바로 부정행위의 직접 증명은 아니다. 사전투표자의 자발적 선택, 정당별 사전투표 독려 전략, 연령·직업·지역 이동성, 코로나19 시기의 투표행태 변화 등 대안 설명을 검토해야 한다. 그러나 모든 지역구에서 같은 방향으로 나타난 2020·2024년의 패턴은 단순 표본오차로 설명하기 어렵고, 독립적 감사 또는 추가 통계분석의 대상이 된다.

이 방향성 문제는 흔히 "대수의 법칙"이라는 말로 설명되지만, 엄밀하게는 "사전투표와 당일투표가 같은 모집단에서 나온 표본이라면 표본오차만으로 얼마나 자주 이런 방향 일치가 생기는가"의 문제이다. 본문에서는 이를 보조 검정으로만 사용하고, 계산의 세부 전제와 결과는 부록 B에 따로 제시한다.

\subsection{6.3 지방선거와의 연결}

2026년 6.3 지방선거에서 보도된 동일 득표쌍 12곳은 위의 사전-당일 득표율 차이 문제와 별개의 현상이다. 그러나 두 현상은 모두 사전투표에서 발생했다는 공통점이 있다.

본 논문이 6.3 지방선거를 중점적으로 보는 이유는 다음과 같다.

\begin{enumerate}
\item 인천 송도1동·송도2동의 동일 득표쌍은 관내사전투표에서 발생했다.
\item 광주전남 10개 지역, 즉 5쌍의 동일 득표쌍도 관내사전투표에서 발생했다.
\item 과거 시·도지사 공식자료에서 한 선거구 안 동일 득표쌍 최대 관측값은 3쌍이었다.
\item 6.3 지방선거 보도 사례는 같은 선거구 안 5쌍으로, 확장 기준선의 과거 최대치를 초과한다.
\item 2020·2024 총선에서도 사전투표는 당일투표와 체계적으로 다른 방향성을 보였다.
\end{enumerate}

따라서 6.3 지방선거의 동일 득표쌍 문제는 고립된 우연 사건으로 보기 어렵다. 2020·2024 총선에서 관찰되는 사전투표의 체계적 편차와 결합하면, "사전투표 계열 자료에서 반복적으로 통계적 이상치가 나타난다"는 더 넓은 문제로 확장된다.

\section{예상 반론과 응답}

본 절은 가장 회의적인 심사자가 제기할 수 있는 반론을 미리 정리하고, 본 연구의 결론이 어느 범위에서 유지되는지를 밝힌다.

\subsection{"개표단위가 많으면 같은 득표쌍은 언제든 나올 수 있다"}

맞다. 본 연구도 동일 득표쌍 한 쌍의 발생을 이상치로 보지 않는다. 개표단위가 많으면 가능한 비교쌍 수가 빠르게 증가하므로, 동일 득표쌍 하나는 자연스럽게 발생할 수 있다. 본 연구의 대상은 "동일 득표쌍의 존재"가 아니라 "같은 선거구, 같은 투표유형, 같은 후보 조합 안에서 5쌍이 반복된 집중도"이다. 확장 기준선에서 과거 시·도지사 선거 51개 선거구 중 한 선거구 안 최대 반복은 3쌍이었다. 따라서 광주전남 5쌍은 단순 중복 발생이 아니라 과거 경험적 상단을 넘는 사건으로 다루어야 한다.

\subsection{"2014년을 넣으면 과거에도 중복이 많다"}

맞다. 2014년 지방선거를 추가하면 시·도지사 동일 득표쌍은 3개에서 15개로 증가하고, 추정 확률도 약 0.0004\%에서 약 0.115\%로 완화된다. 이 점은 본 연구에 불리한 방향의 수정이다. 그럼에도 확장 기준선에서 한 선거구 안 최대 반복은 3쌍이고, 5쌍 이상은 관찰되지 않았다. 따라서 2014년을 포함하는 보수적 기준선에서도 2026년 광주전남 5쌍 반복은 과거 경험분포 밖에 있다.

\subsection{"과거에도 3쌍이 있었다면 5쌍도 큰 차이가 아니다"}

그렇지 않다. 이 주장은 3쌍과 5쌍을 선형적으로 비교하기 때문에 생기는 오해다. 동일 득표쌍 반복은 한 쌍씩 늘어날 때마다 확률이 빠르게 작아지는 꼬리 사건이다. 본 연구의 보수적 기준선에서 3쌍 이상은 약 4.23\%로, 51개 과거 선거구 중 한 번 관찰될 수 있는 수준이다. 반면 5쌍 이상은 약 0.115\%로, 3쌍 이상보다 약 37배 더 드문 사건이다. 따라서 과거 최대값 3쌍은 "동일 득표쌍 반복이 어느 정도까지 자연적으로 관찰되었는가"를 보여주는 기준선이고, 2026년 광주전남 5쌍은 그 기준선을 초과한 사건이다.

또한 본 연구는 3쌍을 숨기거나 제외하지 않았다. 오히려 2014년 자료를 포함해 과거 중복을 더 많이 인정한 뒤에도 5쌍 이상이 없었다는 점을 제시한다. 이는 연구자에게 불리한 자료를 포함한 보수적 검정이다.

\subsection{"읍면동 규모와 지역 성향이 다르므로 단순 확률계산이 틀릴 수 있다"}

맞다. 읍면동 득표쌍은 서로 완전히 독립적이지도 않고 같은 분포를 따르지도 않는다. 그래서 본 연구는 균등분포를 가정해 이론적 범위를 임의로 만든 것이 아니라, 과거 시·도지사 관내사전투표 자료에서 실제 동일쌍 발생 빈도를 역산한 경험적 \(K\)를 사용했다. 또한 \(K=50,000\)처럼 가능한 득표쌍 범위를 더 작게 잡는 민감도 분석도 제시했다. 이 경우에도 5쌍 이상은 약 2.05\%로, 흔한 사건이라고 보기는 어렵다.

\subsection{"전국 12곳은 나중에 찾아낸 사례이므로 확률을 그대로 볼 수 없다"}

맞다. 전국 12곳은 여러 선거와 지역을 훑은 뒤 발견된 사건이므로 나중에 찾아낸 사례라는 문제가 있다. 그래서 본 논문은 전국 6쌍을 핵심 검정으로 삼지 않는다. 핵심 검정은 광주전남 내부 5쌍이다. 같은 후보 조합, 같은 관내사전투표 유형, 같은 광역 선거구 안에 집중된 반복이기 때문에 나중에 찾아낸 사례라는 문제가 훨씬 작다. 전국 12곳은 보조적 맥락이지, 본 논문의 확률 산정 대상이 아니다.

\subsection{"2026년 자료가 아직 공식 통합 파일이 아니라 언론 보도 기반이다"}

이 반론은 초기 초안에는 타당했지만, 현재 원고에는 그대로 적용되지 않는다. 본 연구는 2026년 12개 사건행을 중앙선거관리위원회 선거통계시스템의 개표단위별 개표결과 HTML에서 직접 추출했다. 재현 스크립트 \texttt{scripts/fetch\_nec\_2026\_duplicate\_cases.py}는 선거ID \texttt{0020260603}, 메뉴 \texttt{VCCP08}, 선거종류 \texttt{시·도지사선거}를 POST 조회하고, 관내사전투표 행에서 후보별 득표수를 파싱한다. 그 결과 \texttt{outputs/nec\_2026\_reported\_duplicate\_cases.csv}에는 12개 사건행이, \texttt{outputs/nec\_2026\_reported\_duplicate\_pairs.csv}에는 6개 동일 득표쌍이 생성된다.

따라서 현재 남는 한계는 "언론 보도 기반"이 아니라 "공식 화면 HTML 추출값 기반이며, 아직 통합 XLSX 파일과 개표상황표 원본까지 대조하지는 않았다"는 것이다. 이는 본 연구의 결론을 약화시키기보다 결론의 층위를 분명히 한다. 공식 선거통계시스템 화면값으로도 동일 득표쌍 반복은 재현된다. 다만 특정 부정행위의 법적 확정에는 개표상황표 원본, 1차 분류기 결과, 재확인표 배분 내역, 입력 로그가 추가로 필요하다.

\subsection{"사전투표와 당일투표 차이는 유권자가 스스로 선택한 결과일 수 있다"}

일부는 설명될 수 있다. 그래서 본 논문은 사전-당일 득표율 차이를 독립적인 부정행위 증거로 사용하지 않는다. 이 분석은 보조 검정이다. 다만 2020년과 2024년 국회의원선거에서 분석 가능한 모든 지역구가 같은 방향을 보였다는 사실은 단순 표본오차로 설명하기 어렵다. 유권자 선택 가설이 충분하려면 연령, 직업, 지역 이동성, 정당별 사전투표 독려, 코로나19 시기 효과 등을 반영한 뒤에도 남는 차이가 사라지는지 보여야 한다.

\subsection{"통계적으로 드문 사건은 실제로도 가끔 발생한다"}

맞다. 낮은 확률은 불가능을 뜻하지 않는다. 따라서 본 논문은 낮은 확률만으로 특정 부정행위를 단정하지 않는다. 그러나 선거무결성 영역에서는 낮은 확률의 구조적 이상치가 관찰되면 원자료 감사로 넘어가는 것이 합리적이다. 본 연구의 결론은 처벌이나 무효 판단이 아니라, 개표상황표 원본, 1차 분류기 결과, 재확인표 배분 내역, 전산 입력 로그를 공개해야 한다는 감사 명제이다.

\subsection{"이 주장은 어떻게 반증될 수 있는가"}

본 연구의 주장은 반증 가능하다. 다음 중 하나가 확인되면 결론은 약화되거나 철회되어야 한다.

\begin{enumerate}
\item 2026년 12곳의 동일 득표쌍 보도값이 공식 개표상황표 원본과 다르다.
\item 광주전남 5쌍이 같은 선거구·같은 후보 조합·같은 관내사전투표 유형이 아니었다.
\item 2014·2018·2022년 외의 동일 제도권 시·도지사 관내사전투표 자료에서 한 선거구 안 5쌍 이상 반복 사례가 다수 확인된다.
\item 2026년 5쌍 모두에 대해 1차 분류기 결과, 재확인표 배분, 입력 로그가 독립적으로 일관되며, 동일성이 자연적 개표 과정에서 발생했음이 문서로 확인된다.
\end{enumerate}

이 반증 조건을 명시하는 이유는 본 연구가 정치적 단정이 아니라 검증 가능한 통계 명제임을 분명히 하기 위해서이다.

\section{원자료 감사 판정 기준}

통계적 이상치가 실제 행정·개표 원인과 연결되는지는 공개 화면값만으로 확정할 수 없다. 따라서 본 연구의 실무적 제안은 단순한 의혹 제기가 아니라, 다음 원자료를 같은 사건 정의에 맞춰 대조하자는 것이다.

\begin{table*}[t]
\centering
\scriptsize
\resizebox{\linewidth}{!}{%
\begin{tabular}{p{0.237\linewidth}p{0.237\linewidth}p{0.237\linewidth}p{0.237\linewidth}}
\toprule
원자료 & 확인할 질문 & 우연 설명을 약화시키는 결과 & 우연 설명을 회복시키는 결과 \\
\midrule
개표상황표 원본 & 공식 화면의 12개 사건행 득표값이 원본과 같은가 & 원본도 같은 12개 사건행과 6개 동일쌍을 보인다 & 화면값과 원본값이 달라 동일쌍이 사라진다 \\
1차 분류기 결과 & 최종 동일쌍이 1차 분류 단계부터 존재했는가 & 1차 결과는 달랐는데 재확인·합산 뒤 동일쌍이 생긴다 & 1차부터 최종까지 동일값이 독립적으로 일관된다 \\
재확인표 후보별 배분 & 재확인표가 특정 후보 방향으로 반복 배분됐는가 & 여러 쌍에서 같은 후보 방향의 보정이 반복된다 & 재확인표 배분이 무작위적이고 동일쌍 생성과 무관하다 \\
전산 입력 로그 & 입력·수정 시점과 담당자 흐름이 자연스러운가 & 동일쌍 관련 행에서 비정상적 수정·재입력 패턴이 집중된다 & 입력 로그가 시간순으로 일관되고 수정 흔적이 없다 \\
관내·관외 분리자료 & 동일쌍이 관내사전투표에만 집중되는가 & 동일 후보 조합의 동일쌍이 관내사전투표에 집중된다 & 관내·관외·당일 자료에서 같은 정도의 반복이 나타난다 \\
투표함·개표배치 인수인계 기록 & 사건행 투표지가 어떤 배치와 처리 흐름을 거쳤는가 & 동일쌍 관련 배치가 같은 장비·시간대·심사 흐름에 집중된다 & 배치 흐름이 독립적이고 시간순으로 끊김 없이 설명된다 \\
\bottomrule
\end{tabular}
}
\end{table*}

이 판정표의 목적은 결론을 미리 정하는 것이 아니다. 반대로 어떤 자료가 나오면 본 연구의 결론이 약화되는지도 함께 적는다. 예를 들어 개표상황표 원본에서 12개 사건행의 값이 공식 화면값과 다르거나, 광주전남 5쌍이 같은 후보 조합·같은 관내사전투표 유형이 아닌 것으로 확인되면 본 연구의 핵심 결론은 수정되어야 한다. 반면 원본 자료가 현재 화면값과 일치하고, 1차 결과와 최종 결과 사이의 변화가 동일쌍을 만드는 방향으로 반복된다면, 현재의 통계적 이상치 결론은 원인 조사 단계로 강화된다.

따라서 본 연구가 요구하는 감사는 전면적이고 추상적인 정치 감사가 아니다. 사건행 12개, 특히 광주전남 5쌍에 대해 위 표의 원자료를 제한적으로 공개하면, 우연·표시 오류·재확인표 배분·전산 입력 문제·투표함 또는 배치 흐름 문제 중 어느 설명이 남는지 좁힐 수 있다.

원자료 공개 뒤의 판정 단계는 네 가지로 나눈다. 첫째, 개표상황표 원본이나 공식 통합 파일에서 사건행 값이 달라 동일쌍이 사라지면 본 연구의 핵심 결론은 약화 또는 철회된다. 둘째, 원본값은 일치하지만 1차 분류기 결과부터 최종값까지 모든 후보별 증감이 재확인표·입력 로그·배치 기록으로 완전히 설명되면 통계적 이상치 결론은 유지되지만 행정적으로 설명된 사건으로 낮아진다. 셋째, 원본값이 일치하고 1차 결과와 최종 결과 사이의 변화가 여러 쌍에서 동일쌍을 만드는 방향으로 반복되면 결론은 단순 감사 요구에서 원인조사 단계로 강화된다. 넷째, 원본값·1차 결과·재확인표·입력 로그·인수인계 기록 중 일부가 결손되거나 서로 맞지 않으면 우연 설명은 회복되지 않으며, 결손 자체가 독립 감사의 대상이 된다.

\section{결론}

2026년 지방선거에서 보도된 관내사전투표 동일 득표쌍 반복 현상은 통계적으로 매우 비개연적이다. 특히 광주전남에서 같은 선거구, 같은 후보 조합, 같은 관내사전투표 유형 안에 5개의 동일 득표쌍이 나타났다는 보도는 과거 시·도지사 선거 경험분포와 중복확률 모형 양쪽에서 강한 이상치로 평가된다.

과거 시·도지사 자료 51개 선거구에서 한 선거구 안 동일 득표쌍의 최대 관측값은 3쌍이었다. 반면 광주전남 보도 사건은 5쌍이다. 가능한 득표쌍 범위를 과거 자료에서 추정해 계산하면 이러한 사건의 포아송 근사 확률은 약 0.115\%이다. 별도의 비복원 재표본추출 검정에서도 200,000회 중 5쌍 이상은 0회였고, 3의 법칙 95\% 상한은 약 0.0015\%였다. 이는 단순 우연이라는 설명을 통계적으로 매우 약하게 만든다.

그러나 통계적 비개연성은 곧바로 특정 부정행위의 직접 증거가 아니다. 본 사건은 법적 판단의 결론이 아니라, 우연가설로 설명하기 어려운 통계적 이상치로 규정하는 것이 타당하다. 다만 이 구분은 결론을 약화시키기 위한 표현이 아니다. 가장 회의적인 해석을 적용하더라도 현재 공개자료만으로는 우연 설명이 충분하지 않으며, 선거관리기관이 이를 단순한 우연 또는 착시로 종결하려면 그 설명을 원자료로 입증해야 한다. 선관위 또는 관련 기관이 의혹을 해소하려면 12곳 전체의 개표상황표 원본, 1차 분류기 결과, 재확인표 후보별 배분 내역, 전산 입력 로그, 투표함·개표배치 인수인계 기록을 공개해야 한다.

본 논문의 핵심 결론은 다음 한 문장으로 요약된다.

\begin{quote}
2026년 관내사전투표 동일 득표쌍 12곳, 특히 광주전남 5쌍 반복은 우연이라고 보기에는 지나치게 낮은 확률의 사건이며, 원자료 공개와 독립 검증 없이는 통계적으로 해소되지 않는 비정상 선거자료이다.
\end{quote}

\section{자료 및 코드 가용성}

본 연구의 과거 기준선 분석은 공개된 중앙선거관리위원회 및 공공데이터포털 자료에 기초한다. 분석에 사용한 원자료 파일, 파싱 스크립트, 확률 계산 스크립트, 출력 CSV는 본 제출 패키지 안에 정리했다. 핵심 산출물은 \texttt{README.md}의 재현 절차에 따라 다시 생성할 수 있다.

2026년 지방선거 동일 득표쌍 사건값은 중앙선거관리위원회 선거통계시스템의 개표단위별 개표결과 HTML에서 직접 추출한 값으로 대조했다. 원본 HTML 캐시는 \texttt{data/nec\_2026\_official\_html/}에, 파싱 결과는 \texttt{outputs/nec\_2026\_reported\_duplicate\_cases.csv}와 \texttt{outputs/nec\_2026\_reported\_duplicate\_pairs.csv}에 포함했다. 다만 공공데이터포털의 공식 통합 XLSX 개표자료 또는 개표상황표 원본이 공개되면, 본 연구의 스크립트를 동일하게 적용해 사건값과 반복 수를 추가 재검산해야 한다.

\enlargethispage{3\baselineskip}
\section{연구윤리 및 이해상충}

본 연구는 공개 선거자료와 공개 보도자료만 사용했으며, 개인 식별정보나 비공개 개인정보를 분석하지 않았다. 저자는 특정 후보자, 정당, 선거관리기관의 법적 책임을 단정하지 않는다. 본 연구의 목적은 통계적으로 낮은 확률의 반복 패턴이 발견되었을 때 필요한 원자료 공개와 독립 검증의 범위를 제시하는 것이다.

본 연구의 결론은 "부정행위가 법적으로 확정되었다"가 아니라 "현재 공개자료만으로는 우연 설명이 충분하지 않으며, 공식 원자료 감사가 필요하다"는 검증 명제이다. 이 구분은 연구윤리상 중요하다. 통계적 이상치를 직접 범죄 증거로 과장하지 않고, 동시에 통계적 이상치를 아무 의미 없는 우연으로 축소하지 않기 위해서이다.

\clearpage
\onecolumn
\section*{부록 A. 주요 용어의 직관적 의미}

본 논문에서 사용하는 통계 용어의 직관적 의미는 다음과 같다.

\begin{center}
\small
\begin{tabular}{>{\raggedright\arraybackslash}p{0.18\linewidth}>{\raggedright\arraybackslash}p{0.76\linewidth}}
\toprule
용어 & 직관적 의미 \\
\midrule
귀무가설 & "이 정도 현상은 우연으로도 나올 수 있다"는 기본 가정 \\
대립가설 & "우연으로 보기 어렵고 별도 설명이 필요하다"는 반대 가정 \\
포아송 근사 & 드문 사건이 여러 번 반복될 확률을 계산할 때 쓰는 표준 근사법 \\
가능한 득표쌍 범위 \(K\) & 관내사전투표에서 실제로 나올 수 있는 1·2위 후보 득표수 조합의 실질적 범위 \\
민감도 분석 & \(K\)나 \(N\) 같은 가정을 바꿔도 결론이 유지되는지 확인하는 점검 \\
비복원 재표본추출 & 과거 실제 자료에서 한 번 뽑은 개표단위를 다시 뽑지 않고 표본을 만들어 반복 패턴을 점검하는 방법 \\
3의 법칙 & 어떤 사건이 \(n\)번 중 0회 관찰됐을 때 그 사건의 95\% 상한을 대략 \(3/n\)으로 보는 보수적 경험규칙 \\
꼬리확률 & 평균적인 상황에서 멀리 떨어진 극단적 사건이 나올 확률 \\
대수의 법칙 & 표본 수가 커질수록 무작위 표본의 비율이 모집단 비율에 가까워진다는 원리 \\
부호검정 & 차이의 크기보다 방향이 한쪽으로 얼마나 반복되는지를 보는 비모수 검정 \\
감사 트리거 & 곧바로 범죄나 부정을 확정하는 증거는 아니지만, 원자료 공개와 독립 검증을 요구할 만큼 강한 이상 신호 \\
\bottomrule
\end{tabular}
\end{center}

\vspace{1em}
\section*{부록 B. 사전투표와 대수의 법칙 검정}

사전투표 논의에서 "대수의 법칙"이라는 표현은 자주 쓰이지만, 이 말만으로 부정행위를 증명할 수는 없다. 대수의 법칙은 표본 수가 커질수록 무작위 표본의 비율이 모집단 비율에 가까워진다는 원리이다. 따라서 이 원리를 사전투표에 적용하려면 사전투표자와 당일투표자가 같은 모집단에서 무작위로 나뉘었다는 전제가 필요하다. 실제 선거에서는 사전투표자가 스스로 사전투표를 선택하므로 이 전제는 완전하지 않다.

그럼에도 대수의 법칙 논의가 완전히 무의미한 것은 아니다. 표본 수가 충분히 큰데도 거의 모든 지역구에서 같은 정당 방향의 사전투표 우위가 반복된다면, 이는 단순 표본오차만으로 보기 어렵다. 본 부록은 이 점을 가장 단순한 부호검정으로 확인한다.

각 지역구에서 민주당 후보의 사전투표 양자득표율이 당일투표 양자득표율보다 높은지를 보았다. 귀무가설은 다음과 같이 둔다.

\begin{quote}
사전투표와 당일투표의 차이가 특별히 한 정당 방향으로 치우치지 않는다면, 각 지역구에서 민주당 사전 우위가 나타날 확률은 대략 1/2이다.
\end{quote}

이 귀무가설 아래에서 \(n\)개 지역구 중 \(x\)개 이상이 민주당 사전 우위일 확률은 다음의 한쪽 꼬리확률이다.

\[
P(X \ge x) = \sum_{k=x}^{n} \binom{n}{k}\left(\frac{1}{2}\right)^n
\]

같은 내용을 확률변수 관점으로도 볼 수 있다. 각 지역구 또는 더 작은 개표단위에서

\[
\Delta = p_E - p_D
\]

를 계산하면, \(\Delta\)는 "사전투표 득표율과 당일투표 득표율의 차이"라는 확률변수이다. 단순 표본오차가 주된 원인이라면 \(\Delta\)의 분포는 0 근처에서 좌우로 흔들려야 한다. 반대로 대부분의 지역구에서 \(\Delta>0\)이고 분포의 중심도 0이 아니라 +10\%p 부근으로 이동한다면, 이는 단순한 무작위 오차가 아니라 체계적 방향성을 뜻한다. 이 방향성이 유권자 자기선택 때문인지, 정당별 동원전략 때문인지, 행정·개표 과정의 이상 때문인지는 별도 자료로 구분해야 하지만, "통계적으로 볼 대상이 아니다"라는 반론은 성립하기 어렵다.

공식 총선자료를 직접 파싱해 계산한 결과는 다음과 같다.

\begin{center}
\scriptsize
\resizebox{\linewidth}{!}{%
\begin{tabular}{p{0.158\linewidth}p{0.158\linewidth}p{0.158\linewidth}p{0.158\linewidth}p{0.158\linewidth}p{0.158\linewidth}}
\toprule
선거 & 분석 지역구 수 & 민주당 사전 우위 & 민주당 사전 열위 & 부호검정 한쪽 p값 & 직관적 빈도 \\
\midrule
2016 총선 & 229 & 211 & 18 & \(2.99 \times 10^{-43}\) & 약 \(3.35 \times 10^{42}\)번에 1번 \\
2020 총선 & 236 & 236 & 0 & \(9.06 \times 10^{-72}\) & 약 \(1.10 \times 10^{71}\)번에 1번 \\
2024 총선 & 245 & 245 & 0 & \(1.77 \times 10^{-74}\) & 약 \(5.65 \times 10^{73}\)번에 1번 \\
\bottomrule
\end{tabular}
}
\end{center}

이 계산의 의미는 제한적이지만 분명하다. 사전투표와 당일투표의 차이가 표본오차처럼 위아래로 흔들리는 현상이라면, 2020년 236개 전 지역구와 2024년 245개 전 지역구에서 같은 방향의 민주당 사전 우위가 나타날 확률은 극히 낮다. 따라서 "큰 표본에서 우연히 그렇게 보였을 뿐"이라는 설명은 이 단순 모형 아래에서는 유지되기 어렵다.

다만 이 결과를 "부정행위의 직접 증명"으로 읽으면 안 된다. 부호검정은 유권자 자기선택, 정당별 사전투표 독려, 연령·직업·지역 이동성, 코로나19 시기 효과 같은 대안 설명을 제거하지 못한다. 이 요인들이 실제로 강하게 작동했다면 \(1/2\) 대칭 귀무가설은 약해진다.

따라서 본 부록의 결론은 다음과 같다. 2020년과 2024년 총선의 사전투표 방향 일치는 단순 표본오차로 보기에는 확률적으로 극히 낮다. 그러나 이 결과는 동일 득표쌍 반복 분석을 보강하는 보조 이상성 지표이며, 법적·행정적 결론에는 지역구별 유권자 구성, 사전투표 독려 효과, 관내·관외 분리자료, 개표상황표 원본, 1차 분류기 결과와의 결합 검증이 필요하다.

\vspace{1em}
\section*{참고자료}

\begin{itemize}
\item TV조선. 2026.06.08. "인천·전남 등 전국 12곳서 '동일 득표' 잇따라...'우연의 일치' 해명에도 논란 확산." https://news.tvchosun.com/site/data/html\_dir/2026/06/08/2026060890338.html
\item 뉴스1/다음. 2026.06.06. "송도1동과 송도2동 표 수가 왜 똑같아?...박찬대·유정복 득표 논란." https://v.daum.net/v/20260606113452062
\item 한겨레. 2026.06.08. "'송도1·2동, 박찬대·유정복 같은 득표수'에 선관위 '부정개표 없어'." https://www.hani.co.kr/arti/area/capital/1262526.html
\item 공공데이터포털. 중앙선거관리위원회\_제6회 전국동시지방선거 개표결과\_20140604. https://www.data.go.kr/data/15048207/fileData.do
\item 공공데이터포털. 중앙선거관리위원회\_투·개표 정보. https://www.data.go.kr/data/15000900/openapi.do
\item 중앙선거관리위원회 국가선거정보 개방포털. 국회의원·대통령·전국동시지방선거 개표결과 공개자료. https://data.nec.go.kr/
\item 공공데이터포털. 중앙선거관리위원회 선거 개표결과 파일데이터. https://www.data.go.kr/
\item 공공데이터포털. 광주광역시\_주민센터 현황. https://www.data.go.kr/data/3082663/fileData.do
\item 전라남도청. 전남소개 일반현황 행정구역. https://www.jeonnam.go.kr/contentsView.do?menuId=jeonnam0602030000
\end{itemize}

\section*{재현 산출물}

\begin{itemize}
\item 의존성 파일: \texttt{requirements.txt}
\item 전체 재현 스크립트: \texttt{scripts/run\_all.py}
\item 분석 스크립트: \texttt{scripts/analyze\_duplicates.py}
\item 시·도지사 실제 1·2위 기준선 스크립트: \texttt{scripts/analyze\_governor\_actual\_top2.py}
\item 시·도지사 비복원 재표본추출 스크립트: \texttt{scripts/bootstrap\_governor\_duplicates.py}
\item 확률 및 민감도 표 재현 스크립트: \texttt{scripts/probability\_sensitivity.py}
\item 사전-당일 검산 스크립트: \texttt{scripts/analyze\_early\_day\_assembly.py}
\item 추가 원자료: \texttt{data/local2014.xlsx}
\item 기존 요약표: \texttt{outputs/duplicate\_summary.csv}
\item 중복 그룹: \texttt{outputs/duplicate\_groups.csv}
\item 중복 사례: \texttt{outputs/duplicate\_examples.csv}
\item 시·도지사 실제 1·2위 기준선 요약: \texttt{outputs/governor\_actual\_top2\_summary.csv}
\item 시·도지사 선거구별 기준선: \texttt{outputs/governor\_actual\_top2\_by\_contest.csv}
\item 시·도지사 동일 득표쌍 사례: \texttt{outputs/governor\_actual\_top2\_duplicates.csv}
\item 시·도지사 비복원 재표본추출 요약: \texttt{outputs/governor\_bootstrap\_summary.csv}
\item 시·도지사 비복원 재표본추출 히스토그램: \texttt{outputs/governor\_bootstrap\_histogram.csv}
\item 핵심 확률표: \texttt{outputs/probability\_core.csv}
\item 가능한 득표쌍 범위 민감도 표: \texttt{outputs/probability\_k\_sensitivity.csv}
\item 개표단위 수 민감도 표: \texttt{outputs/probability\_n\_sensitivity.csv}
\item 사전-당일 검산 요약: \texttt{outputs/early\_day\_assembly\_summary.csv}
\item 사전-당일 지역구별 결과: \texttt{outputs/early\_day\_assembly\_twoparty.csv}
\item 본 논문: \texttt{paper\_statistical\_implausibility\_ko.md}
\item IEIE LaTeX 원고: \texttt{latex/ieie/main.tex}
\item IEIE PDF 원고: \texttt{latex/ieie/main.pdf}
\end{itemize}

\end{document}
